% =====================================================================
% RELATÓRIO 11 — POO com Java
% Capítulo para ser incorporado à Apostila S01 (OO + PF), após o Relatório 10.
% =====================================================================

% ---------------------------------------------------------------------
% Definições usadas neste capítulo (podem ser movidas para o preâmbulo).
% - java: estilo que usa a linguagem Java do listings e corrige a
%   acentuação dentro dos códigos (mesma solução dos estilos anteriores).
% - azulRelatorio e exercicio: já são criados no Relatório 7. As linhas
%   abaixo só os recriam se este capítulo for compilado sem ele.
% ---------------------------------------------------------------------
\lstdefinestyle{java}{
    language=Java,
    literate=
        {á}{{\'a}}1 {à}{{\`a}}1 {ã}{{\~a}}1 {â}{{\^a}}1
        {é}{{\'e}}1 {ê}{{\^e}}1 {è}{{\`e}}1 {í}{{\'\i}}1
        {ó}{{\'o}}1 {õ}{{\~o}}1 {ô}{{\^o}}1 {ú}{{\'u}}1
        {ç}{{\c{c}}}1 {Á}{{\'A}}1 {É}{{\'E}}1 {Ç}{{\c{C}}}1
        {Ã}{{\~A}}1 {Â}{{\^A}}1 {Ê}{{\^E}}1 {Õ}{{\~O}}1
        {—}{{\textemdash}}1 {Í}{{\'I}}1 {Ó}{{\'O}}1 {Ú}{{\'U}}1
}

\definecolor{azulRelatorio}{rgb}{0.1,0.3,0.8}

\ifcsname exercicio\endcsname\else
\newtcolorbox{exercicio}[1]{
    colback=white,
    colframe=azulRelatorio,
    coltitle=white,
    fonttitle=\normalfont,
    boxrule=0.8pt,
    arc=1.5mm,
    title={#1}
}
\fi

\chapter{Relatório 11: Programação Orientada a Objetos em Java}

\section{Fundamentos Teóricos e Pilares da POO}

No Relatório 10, o JavaScript mostrou uma linguagem de tipagem dinâmica, em que os contratos
entre classes só são conferidos durante a execução. O \textbf{Java} vai no sentido oposto. Ele
tem \textbf{tipagem estática}: todo atributo, parâmetro, variável e retorno de método declara o
seu tipo. O código é \textbf{compilado} antes de rodar, e o resultado roda na \textbf{JVM}
(\textit{Java Virtual Machine}), o que permite que o mesmo programa funcione em qualquer sistema
operacional. Em Java, \textbf{tudo fica dentro de uma classe}: não existem funções soltas, e até
o programa principal é um método. As interfaces e as classes abstratas são nativas, e o
compilador confere os contratos antes de o programa rodar. A sintaxe lembra muito a do C\#
(Relatório 7): chaves, ponto e vírgula e modificadores de acesso como \texttt{public} e
\texttt{private}. Assim como nas linguagens anteriores, a memória é liberada automaticamente
pelo \textbf{Garbage Collector}.

A tabela abaixo relaciona as estruturas já vistas em JavaScript com as equivalentes em Java:

\begin{center}
\small
\begin{tabular}{|l|l|l|}
\hline
\textbf{Conceito} & \textbf{JavaScript (Relatório 10)} & \textbf{Java (Relatório 11)} \\
\hline
Tipos & não são declarados & declarados: \texttt{String nome} \\
Construtor & \texttt{constructor(nome)} & \texttt{public Autor(String nome)} \\
Instanciação & \texttt{new Autor(nome)} & \texttt{Autor a = new Autor(nome);} \\
Atributo privado & \texttt{\#nome} & \texttt{private String nome;} \\
Getter & \texttt{get nome() \{ ... \}} & \texttt{public String getNome()} \\
Herança & \texttt{class Cliente extends Pessoa} & igual, com \texttt{extends} \\
Construtor do pai & \texttt{super(nome)} & \texttt{super(nome);} \\
Sobrescrita & basta redefinir o método & redefinir + \texttt{@Override} \\
Interface & \texttt{throw new Error()} na base & \texttt{interface} + \texttt{implements} \\
Classe abstrata & não existe & \texttt{abstract class} \\
Lista & \texttt{[]} e \texttt{.push()} & \texttt{ArrayList<Tipo>} e \texttt{.add()} \\
Dicionário & \texttt{new Map()} e \texttt{.set()} & \texttt{HashMap<Chave, Valor>} e \texttt{.put()} \\
Repetição & \texttt{lista.forEach(x => ...)} & \texttt{for (Tipo x : lista)} \\
Saída & \texttt{console.log(...)} & \texttt{System.out.println(...)} \\
\hline
\end{tabular}
\end{center}

Os quatro pilares continuam valendo, agora ilustrados com o sistema de uma biblioteca, com
autores, livros, estantes, clientes e bibliotecários:

\begin{itemize}
    \item \textbf{Abstração}: Cada objeto guarda apenas o que importa para a biblioteca. Um
    livro precisa de ISBN, título e autor.

    \item \textbf{Encapsulamento}: Os atributos de \texttt{Autor}, \texttt{Livros} e
    \texttt{Estante} são \texttt{private} e só são lidos pelos getters. Os de \texttt{Pessoa} são
    \texttt{protected}.

    \item \textbf{Herança}: \texttt{Cliente} e \texttt{Bibliotecario} reaproveitam a estrutura de
    \texttt{Pessoa} (relação ``\textbf{É UM}'': um Cliente é uma Pessoa).

    \item \textbf{Polimorfismo}: Um \texttt{Cliente} ou um \texttt{Bibliotecario} pode ser
    guardado em uma variável do tipo \texttt{Pessoa}, e qualquer classe que implemente
    \texttt{GerenciaInventario} pode ser usada onde a interface é esperada. Quando a filha
    sobrescreve um método, é a versão dela que roda.
\end{itemize}

\subsection{Interfaces e Classes Abstratas: Contratos Conferidos pelo Compilador}

A novidade deste relatório é que o Java tem as duas formas de contrato \textbf{nativas}:

\begin{itemize}
    \item \textbf{Interface} (\texttt{interface}): Define apenas \textbf{quais} métodos uma
    classe deve ter, sem implementação. A classe que assina o contrato usa \texttt{implements}
    e é obrigada a implementar todos os métodos. No código de referência, \texttt{Bibliotecario}
    implementa a interface \texttt{GerenciaInventario}.

    \item \textbf{Classe abstrata} (\texttt{abstract class}): Fica no meio do caminho. Ela pode
    ter atributos, construtor e métodos prontos, e também \textbf{métodos abstratos}, que só têm
    a assinatura. A classe filha usa \texttt{extends} e é obrigada a implementar os métodos
    abstratos. Uma classe abstrata nunca pode ser instanciada. No código de referência,
    \texttt{Pessoa} é abstrata.
\end{itemize}

Use uma interface quando classes sem nenhuma relação entre si precisam cumprir o mesmo contrato,
e uma classe abstrata quando as filhas compartilham uma base comum, mas cada uma precisa
completar uma parte. Uma classe herda de uma única classe, mas pode implementar várias
interfaces.

A grande diferença em relação ao Ruby e ao JavaScript é o \textbf{momento} em que o contrato é
conferido. Naquelas linguagens, o erro só aparecia quando o método que faltava era chamado,
durante a execução. Em Java, o \textbf{compilador} confere tudo antes: se uma classe esquecer
um método obrigatório, ou se o programa tentar instanciar uma classe abstrata, o código nem
chega a compilar.

\subsection{Relacionamentos entre Classes}

A composição e a agregação, vistas nos relatórios anteriores, continuam valendo (relação
``\textbf{TEM UM}''):

\begin{itemize}
    \item \textbf{Composição (relação forte)}: A parte é criada \textbf{dentro} do objeto dono e
    ``morre'' junto com ele. Exemplo: o \texttt{HashMap} do inventário é criado dentro do
    construtor da própria \texttt{Estante}.

    \item \textbf{Agregação (relação fraca)}: A parte é criada \textbf{fora} e o dono apenas
    guarda uma referência a ela. Exemplo: o \texttt{Autor} é criado antes e só depois é entregue
    ao livro, e os livros são criados no programa principal e depois entregues à estante e ao
    cliente.
\end{itemize}

\section{Documentação Sintática das Estruturas de POO}

Esta seção detalha todas as palavras-chave e construções sintáticas utilizadas na implementação
de POO em Java. Os exercícios deste relatório devem ser resolvidos \textbf{apenas} com as
estruturas documentadas aqui.

\subsection{Estrutura do Programa, Tipos e Saída no Console}

Todo programa Java começa pelo método \texttt{main}, que fica dentro de uma classe. Comentários
de linha começam com \texttt{//}, os blocos ficam entre chaves e as instruções terminam com
ponto e vírgula. Toda variável é declarada com o seu \textbf{tipo} antes do nome:

\begin{lstlisting}[style=java]
public class Main {
    public static void main(String[] args) {
        // comentários de linha começam com //
        String titulo = "Neuromancer";   // texto
        int ano = 1984;                  // número inteiro
        double nota = 9.5;               // número com casas decimais
        boolean disponivel = true;       // verdadeiro ou falso

        System.out.println("Livro: " + titulo + " (" + ano + ")");
    }
}
\end{lstlisting}

\begin{itemize}
    \item \texttt{public static void main(String[] args)}: Início do programa (o equivalente
    à \texttt{Main} do C\#). Essa linha é sempre escrita exatamente assim.

    \item \texttt{System.out.println(...)}: Imprime um valor e pula uma linha. O \texttt{+} junta
    textos, e números são convertidos em texto automaticamente.

    \item \textbf{Tipos básicos}: \texttt{String} (texto, entre aspas duplas), \texttt{int}
    (inteiro), \texttt{double} (número com casas decimais, escrito com ponto, como
    \texttt{9.5}) e \texttt{boolean} (\texttt{true} ou \texttt{false}).

    \item \texttt{if (condição) \{ ... \} else \{ ... \}}: Estrutura de decisão. As comparações
    são \texttt{==}, \texttt{!=}, \texttt{>}, \texttt{<}, \texttt{>=} e \texttt{<=}, e as
    condições podem ser combinadas com \texttt{\&\&} (e), \texttt{||} (ou) e \texttt{!} (não).

    \item \textbf{Comparar textos}: Para saber se duas \texttt{String} têm o mesmo conteúdo, use
    \texttt{texto.equals(outro)} (ou \texttt{equalsIgnoreCase}, que ignora maiúsculas e
    minúsculas). O \texttt{==} entre textos não compara o conteúdo.
\end{itemize}

\subsection{Um Arquivo por Exercício}

O código de referência é um projeto completo: cada classe pública fica em um arquivo com o
mesmo nome da classe (\texttt{Pessoa.java}, \texttt{Cliente.java}...), e todos os arquivos
começam com \texttt{package org.example;}. Nos exercícios, para simplificar, todas as classes
de um exercício podem ficar em um \textbf{único arquivo}:

\begin{lstlisting}[style=java]
import java.util.ArrayList;   // os imports ficam no topo do arquivo

public class Main {           // a classe com o main vem primeiro
    public static void main(String[] args) {
        Ajudante ajudante = new Ajudante();
        ajudante.ajudar();
    }
}

class Ajudante {              // as outras classes vêm depois, sem public
    public void ajudar() {
        System.out.println("Ajudando!");
    }
}
\end{lstlisting}

\begin{itemize}
    \item \textbf{Execução}: No terminal, o comando \texttt{java Main.java} compila e executa o
    arquivo de uma vez (a partir do Java 11). A primeira classe do arquivo precisa ter o
    \texttt{main}.

    \item \texttt{public}: Só uma classe do arquivo pode ser \texttt{public}: a do
    \texttt{main}. As demais são declaradas só com \texttt{class}.

    \item \textbf{Erros de compilação}: Se o código tiver um erro que o compilador detecta, o
    Java mostra a linha do erro e a mensagem \texttt{error: compilation failed}, e
    \textbf{nenhuma} linha do programa é executada.
\end{itemize}

\subsection{Classes, Atributos e Construtores}

Os atributos são declarados no corpo da classe, com o modificador de acesso, o tipo e o nome. O
\textbf{construtor} tem o mesmo nome da classe e não tem tipo de retorno. Os métodos declaram o
tipo do valor que retornam:

\begin{lstlisting}[style=java]
public class Autor {

    private String nome;               // modificador + tipo + nome

    public Autor(String nome) {        // construtor: mesmo nome da classe
        this.nome = nome;
    }

    public String getNome() {          // este método retorna uma String
        return nome;
    }
}

Autor kazuo = new Autor("Kazuo Ishiguro");   // Tipo variavel = new ...
System.out.println(kazuo.getNome());
\end{lstlisting}

\begin{itemize}
    \item \textbf{Construtor}: Executado automaticamente pelo \texttt{new}. Os parâmetros também
    declaram tipo: \texttt{public Livros(String isbn, String titulo, Autor autor)}.

    \item \texttt{this}: Referencia o próprio objeto (como no C\# e no JavaScript). É usado para
    diferenciar o atributo \texttt{this.nome} do parâmetro \texttt{nome}.

    \item \textbf{Tipo de retorno}: Vem antes do nome do método. \texttt{void} indica que o
    método não retorna nada; \texttt{String}, \texttt{double} ou o nome de uma classe indicam o
    tipo do valor devolvido com \texttt{return}.

    \item \textbf{Declaração com} \texttt{new}: A variável que guarda o objeto também declara
    tipo: \texttt{Autor kazuo = new Autor(...);}.
\end{itemize}

\subsection{Modificadores de Acesso, Getters e Setters}

O nível de acesso de cada atributo e método é indicado por um modificador, como no C\#:

\begin{lstlisting}[style=java]
public abstract class Pessoa {

    protected String nome;   // acessível nas filhas e no mesmo pacote
    protected int id;

    public Pessoa(String nome, int id) {
        this.nome = nome;
        this.id = id;
    }

    public String getNome() {           // getter
        return nome;
    }

    public void setNome(String nome) {  // setter
        this.nome = nome;
    }
}
\end{lstlisting}

\begin{itemize}
    \item \texttt{public}: Acessível de qualquer classe.

    \item \texttt{private}: Acessível só dentro da própria classe. Nem as classes filhas
    enxergam um atributo \texttt{private} do pai; para ler o valor, elas usam o getter, como o
    \texttt{Cliente} faz com \texttt{getNome()}. Acessar um atributo \texttt{private} fora da
    classe é um erro de compilação (\texttt{nome has private access}).

    \item \texttt{protected}: Acessível nas classes filhas e também em todas as classes do mesmo
    pacote.

    \item \textbf{Getter e setter}: Métodos públicos que leem (\texttt{getNome}) e alteram
    (\texttt{setNome}) um atributo de forma controlada. Por convenção, os nomes começam com
    \texttt{get} e \texttt{set}.
\end{itemize}

\subsection{Herança com \texttt{extends} e \texttt{super(...)}}

A classe pai é indicada com \texttt{extends}. O construtor da filha chama \texttt{super(...)}
para executar o construtor do pai, repassando os parâmetros de que ele precisa:

\begin{lstlisting}[style=java]
public class Cliente extends Pessoa {   // Cliente É UMA Pessoa

    private List<Livros> livrosEmprestados;

    public Cliente(String nome, int id) {
        super(nome, id);   // chama o construtor de Pessoa
        this.livrosEmprestados = new ArrayList<>();
    }
}
\end{lstlisting}

\begin{itemize}
    \item \texttt{class Filha extends Pai}: Declara a herança. Uma classe tem um único pai.

    \item \texttt{super(...)}: Deve ser a \textbf{primeira} instrução do construtor da filha (o
    equivalente ao \texttt{super(...)} do JavaScript e do Ruby e ao \texttt{: base(...)} do C\#).

    \item \textbf{Construtor obrigatório na filha}: Se o construtor do pai recebe parâmetros, a
    filha precisa declarar o seu próprio construtor e chamar \texttt{super(...)} nele, mesmo que
    ela não tenha nenhum atributo novo.
\end{itemize}

Uma variável declarada com o tipo da classe pai pode guardar um objeto de qualquer classe
filha. Pela variável, só podem ser chamados os métodos que existem no pai, mas, se a filha
sobrescreveu algum deles, é a versão da filha que roda. O exemplo abaixo é complementar e não
faz parte do código de referência:

\begin{lstlisting}[style=java]
Pessoa pessoa = new Cliente("Gerard Way", 101);   // pai guardando filha
System.out.println(pessoa.getNome());              // método de Pessoa
\end{lstlisting}

\subsection{Classes e Métodos Abstratos}

Uma classe abstrata é declarada com \texttt{abstract class}. Um \textbf{método abstrato} também
leva a palavra \texttt{abstract} e tem só a assinatura, terminada em ponto e vírgula, sem corpo.
A classe \texttt{Pessoa} do código de referência é abstrata, mas não tem métodos abstratos; o
exemplo abaixo, complementar, mostra como eles funcionam:

\begin{lstlisting}[style=java]
public abstract class Pessoa {
    // ... atributos, construtor e getters ...

    // método abstrato: só a assinatura, sem corpo
    public abstract void apresentar();
}

public class Cliente extends Pessoa {
    // ... construtor ...

    @Override
    public void apresentar() {   // a filha é obrigada a implementar
        System.out.println("Sou o cliente " + getNome());
    }
}

Pessoa alguem = new Pessoa("Alguém", 1);
// erro de compilação: Pessoa is abstract; cannot be instantiated
\end{lstlisting}

\begin{itemize}
    \item \textbf{Instanciação proibida}: Tentar criar um objeto de uma classe abstrata com
    \texttt{new} é um erro de compilação, e o programa nem chega a rodar.

    \item \textbf{Implementação obrigatória}: Toda filha (que não seja abstrata também) precisa
    implementar todos os métodos abstratos do pai. Se esquecer algum, o código não compila.

    \item \textbf{Partes prontas}: Além dos métodos abstratos, a classe abstrata pode ter
    atributos, construtor e métodos comuns, que as filhas herdam normalmente.
\end{itemize}

\subsection{Sobrescrita com \texttt{@Override}}

A filha sobrescreve um método escrevendo outro com a \textbf{mesma assinatura} (mesmo nome, mesmos
parâmetros e mesmo tipo de retorno). A anotação \texttt{@Override}, escrita logo acima,
pede ao compilador que confira se o método realmente existe no pai ou na interface; se não
existir (por exemplo, por um erro de digitação), o código não compila. No código de
referência, ela aparece na implementação de um método de interface:

\begin{lstlisting}[style=java]
@Override
public Livros localizarLivro(String titulo, Estante estante) {
    // ...
}
\end{lstlisting}

O mesmo vale para sobrescrever um método comum, que já tem corpo no pai. O exemplo abaixo é
complementar e não faz parte do código de referência:

\begin{lstlisting}[style=java]
// Na Classe Base (Pessoa): método comum, com corpo
public String descricao() {
    return "Pessoa: " + nome;
}

// Na Classe Derivada (Cliente): sobrescreve o método do pai
@Override
public String descricao() {
    return "Cliente: " + getNome();
}
\end{lstlisting}

Dentro do método sobrescrito, \texttt{super.descricao()} chama a versão do pai (o equivalente ao
\texttt{base.Metodo()} do C\#).

\subsection{Interfaces com \texttt{interface} e \texttt{implements}}

Uma interface é declarada com \texttt{interface} e contém só as assinaturas dos métodos. A
classe que a implementa usa \texttt{implements} e escreve o corpo de cada método:

\begin{lstlisting}[style=java]
public interface GerenciaInventario {
    // só a assinatura: quem implementar decide como funciona
    Livros localizarLivro(String titulo, Estante estante);
}

public class Bibliotecario extends Pessoa implements GerenciaInventario {
    // ... construtor ...

    @Override
    public Livros localizarLivro(String titulo, Estante estante) {
        // implementação do contrato
    }
}
\end{lstlisting}

\begin{itemize}
    \item \textbf{Métodos da interface}: São escritos sem corpo e sem \texttt{public}, mas a
    classe que os implementa deve declará-los como \texttt{public}.

    \item \texttt{implements}: Uma classe pode implementar várias interfaces, separadas por
    vírgula, e ao mesmo tempo herdar de uma classe com \texttt{extends}.

    \item \textbf{Interface como tipo}: Uma variável, um parâmetro ou uma coleção podem ter o
    tipo da interface. Assim, eles aceitam objetos de qualquer classe que a implemente.
\end{itemize}

\subsection{Composição e Agregação}

A diferença entre os dois relacionamentos está em \textbf{onde a parte é criada}:

\begin{lstlisting}[style=java]
// COMPOSIÇÃO: o HashMap é criado DENTRO do construtor da Estante
this.inventario = new HashMap<>();

// AGREGAÇÃO: o Autor é criado FORA e só recebido pelo Livro
public Livros(String isbn, String titulo, Autor autor) {
    this.autor = autor;
}
\end{lstlisting}

Com uma classe própria, a composição funciona do mesmo jeito: basta criar a parte com
\texttt{new} dentro do construtor do dono. O exemplo abaixo é complementar e não faz parte do
código de referência:

\begin{lstlisting}[style=java]
public class Biblioteca {
    private Estante estante;

    public Biblioteca(String codigoDaEstante) {
        // composição: a Estante nasce junto com a Biblioteca
        this.estante = new Estante(codigoDaEstante);
    }
}
\end{lstlisting}

\subsection{\texttt{ArrayList} e for-each}

O \texttt{ArrayList} é a lista de tamanho dinâmico do Java. Entre os sinais de menor e maior,
indica-se o tipo dos elementos; uma lista do tipo da classe pai aceita objetos de qualquer
classe filha:

\begin{lstlisting}[style=java]
import java.util.ArrayList;   // no topo do arquivo

ArrayList<Livros> emprestados = new ArrayList<>();   // lista vazia
emprestados.add(billsBook);                          // adiciona no final
emprestados.add(neuromancer);
System.out.println(emprestados.size());              // quantidade
System.out.println(emprestados.get(0).getTitulo());  // posição 0

for (Livros livro : emprestados) {                   // for-each
    System.out.println("- " + livro.getTitulo());
}
\end{lstlisting}

\begin{itemize}
    \item \texttt{ArrayList<Tipo>}: Lista que guarda objetos do tipo informado. O
    \texttt{import java.util.ArrayList;} precisa estar no topo do arquivo.

    \item \texttt{.add(objeto)}, \texttt{.size()}, \texttt{.get(posicao)} e
    \texttt{.contains(objeto)}: Inserem no final, retornam a quantidade, retornam o elemento da
    posição (a primeira é 0) e testam se o objeto está na lista.

    \item \texttt{for (Tipo variavel : lista)}: O \textbf{for-each}, que percorre a lista
    elemento por elemento (o equivalente ao \texttt{forEach} do JavaScript e ao \texttt{for x in
    lista} do Python).

    \item \texttt{List<Tipo>}: No código de referência, alguns atributos são declarados como
    \texttt{List} e criados como \texttt{ArrayList}. \texttt{List} é a interface que o
    \texttt{ArrayList} implementa; para os exercícios, \texttt{ArrayList} basta.
\end{itemize}

\subsection{\texttt{HashMap}, \texttt{HashSet} e \texttt{null}}

O \texttt{HashMap} guarda pares \textbf{chave e valor}, e cada chave é única. Entre os sinais de
menor e maior, indica-se o tipo da chave e o tipo do valor. Os trechos com valores fixos abaixo
são complementares:

\begin{lstlisting}[style=java]
import java.util.HashMap;   // no topo do arquivo

HashMap<String, Livros> inventario = new HashMap<>();
inventario.put("978-8576572022", neuromancer);     // chave -> valor

Livros achado = inventario.get("978-8576572022");  // devolve o livro
Livros nada = inventario.get("000");               // chave não existe

if (nada == null) {
    System.out.println("Nenhum livro com esse ISBN.");
}

for (Livros livro : inventario.values()) {   // percorre só os valores
    System.out.println(livro.getTitulo());
}
\end{lstlisting}

\begin{itemize}
    \item \texttt{.put(chave, valor)}: Guarda o par. Se a chave já existir, o valor antigo é
    substituído.

    \item \texttt{.get(chave)}: Devolve o valor guardado com aquela chave, ou \texttt{null} se a
    chave não existir.

    \item \texttt{null}: Valor que significa ``nenhum objeto''. Ele é testado com
    \texttt{== null} ou \texttt{!= null}, como no programa principal do código de referência.

    \item \texttt{.containsKey(chave)} e \texttt{.values()}: Testam se a chave existe e
    devolvem só os valores, para serem percorridos com o for-each.

    \item \texttt{HashSet<Tipo>}: Coleção que \textbf{não aceita elementos repetidos} e não
    guarda a ordem de inserção. Seu \texttt{.add(objeto)} devolve \texttt{false} quando o
    elemento já estava lá. Ela não é necessária nos exercícios.
\end{itemize}

\section{Código de Referência Completo}

Abaixo está a implementação de referência utilizada para exemplificação do paradigma POO em
Java. O projeto tem um arquivo por classe, todos no pacote \texttt{org.example}:

\begin{lstlisting}[style=java, caption={Exemplo de Aplicação de POO em Java: \texttt{GerenciaInventario.java}}]
package org.example;

// INTERFACE
// define um contrato — qualquer classe que implementar GerenciaInventario
// É OBRIGADA a ter o método localizarLivro (o compilador confere isso)
// equivale ao module com raise NotImplementedError do Ruby ou ao ABC do
// Python

public interface GerenciaInventario {
    // o método é declarado mas não implementado aqui
    // quem implementar a interface decide como ele vai funcionar
    Livros localizarLivro(String titulo, Estante estante);
}
\end{lstlisting}

\begin{lstlisting}[style=java, caption={Exemplo de Aplicação de POO em Java: \texttt{Pessoa.java}}]
package org.example;

// CLASSE ABSTRATA
// Pessoa é o molde principal — não pode ser instanciada diretamente
// só existe para ser herdada por Cliente, Bibliotecario, Gerente etc
// abstract em Java equivale ao ABC do Python
public abstract class Pessoa {

    // ENCAPSULAMENTO
    // protected libera o acesso direto para as classes filhas e também
    // para todas as classes do mesmo pacote (aqui, org.example)
    // é parecido com o _ (underline) do Python e com o protected do C#
    protected String nome;
    protected int id;

    public Pessoa(String nome, int id) {
        this.nome = nome;
        this.id = id;
    }

    // GETTER e SETTER
    // como os atributos são protected, criamos métodos públicos
    // pra ler e alterar de forma controlada de fora da hierarquia
    public String getNome() {
        return nome;
    }

    public int getId() {
        return id;
    }

    public void setNome(String nome) {
        this.nome = nome;
    }
}
\end{lstlisting}

\begin{lstlisting}[style=java, caption={Exemplo de Aplicação de POO em Java: \texttt{Autor.java}}]
package org.example;

import java.util.ArrayList;
import java.util.List;

// CLASSE SIMPLES
// Autor não herda de nada, é só um molde com nome e títulos escritos
// vai ser referenciado dentro de Livros (isso é Agregação)
public class Autor {

    // ENCAPSULAMENTO
    // private significa que só a própria classe pode acessar diretamente
    // o acesso de fora é feito pelos getters e setters
    private String nome;
    private List<String> titulosEscritos;

    public Autor(String nome) {
        this.nome = nome;
        // ArrayList é a lista mais comum do Java
        // equivale ao List do C# ou ao array do Ruby
        this.titulosEscritos = new ArrayList<>();
    }

    // GETTER
    // como nome é private, criamos um método público pra lê-lo de fora
    public String getNome() {
        return nome;
    }

    public void adicionarTitulo(String titulo) {
        if (!titulosEscritos.contains(titulo)) {
            titulosEscritos.add(titulo);
            System.out.println("[Autor] Título '" + titulo
                    + "' adicionado com sucesso.");
        } else {
            System.out.println("[Autor] Aviso: Título '" + titulo
                    + "' já existe.");
        }
    }

    public List<String> getTitulosEscritos() {
        return titulosEscritos;
    }
}
\end{lstlisting}

\begin{lstlisting}[style=java, caption={Exemplo de Aplicação de POO em Java: \texttt{Livros.java}}]
package org.example;

import java.util.ArrayList;

// AGREGAÇÃO
// Livros guarda uma referência a um Autor que já existia antes
// se o Livro sumir, o Autor continua existindo
public class Livros {

    private String isbn;
    private String titulo;
    private Autor autor;
    private ArrayList<String> palavrasChave;

    public Livros(String isbn, String titulo, Autor autor) {
        this.isbn = isbn;
        this.titulo = titulo;
        this.autor = autor;
        this.palavrasChave = new ArrayList<>();

        // ao criar o Livro, já notificamos o Autor que ele escreveu esse
        // título
        // o Livro conhece o Autor e usa os métodos dele, mas não é dono
        // dele: o Autor foi criado fora e continua existindo sem o Livro
        autor.adicionarTitulo(titulo);
    }

    public String getIsbn() {
        return isbn;
    }

    public String getTitulo() {
        return titulo;
    }

    public void adicionarPalavraChave(String palavra) {
        this.palavrasChave.add(palavra);
    }

    public String getDetalhes() {
        return "'" + titulo + "' (ISBN: " + isbn + ") - Autor: "
                + autor.getNome();
    }
}
\end{lstlisting}

\begin{lstlisting}[style=java, caption={Exemplo de Aplicação de POO em Java: \texttt{Estante.java}}]
package org.example;

import java.util.HashMap;

// AGREGAÇÃO
// a Estante guarda Livros que já existiam fora dela (criados no main)
// ela mapeia os Livros pelo ISBN usando um HashMap
public class Estante {

    private String codigoEstante;

    // HashMap é uma estrutura de chave -> valor
    // aqui a chave é o ISBN (String) e o valor é o objeto Livros
    // equivale ao dicionário do Python ou ao Hash do Ruby
    private HashMap<String, Livros> inventario;

    public Estante(String codigo) {
        this.codigoEstante = codigo;
        // COMPOSIÇÃO: o HashMap é criado aqui dentro e só existe na Estante
        this.inventario = new HashMap<>();
    }

    public void adicionarLivro(Livros livro) {
        // put(chave, valor) adiciona um par no HashMap
        this.inventario.put(livro.getIsbn(), livro);
        System.out.println("[Inventário] Livro adicionado à estante "
                + codigoEstante + ": " + livro.getTitulo());
    }

    public Livros buscarPorIsbn(String isbn) {
        // get(chave) busca um valor pelo ISBN no HashMap
        return this.inventario.get(isbn);
    }

    public HashMap<String, Livros> getInventario() {
        return inventario;
    }

    public String getCodigoEstante() {
        return codigoEstante;
    }
}
\end{lstlisting}

\begin{lstlisting}[style=java, caption={Exemplo de Aplicação de POO em Java: \texttt{Cliente.java}}]
package org.example;

import java.util.ArrayList;
import java.util.List;

// HERANÇA
// Cliente herda de Pessoa com extends
// a relação é: "Cliente É UMA Pessoa"
public class Cliente extends Pessoa {

    // AGREGAÇÃO
    // livrosEmprestados guarda referências pra Livros que já existiam
    // fora do Cliente
    // se o Cliente sumir, os Livros continuam existindo
    private List<Livros> livrosEmprestados;

    public Cliente(String nome, int id) {
        // super() chama o construtor da classe mãe (Pessoa)
        // obrigatório em Java quando a classe pai tem construtor com
        // parâmetros
        super(nome, id);
        this.livrosEmprestados = new ArrayList<>();
    }

    public void emprestarLivro(Livros livro) {
        if (!livrosEmprestados.contains(livro)) {
            livrosEmprestados.add(livro);
            System.out.println("[Empréstimo] Cliente " + getNome()
                    + " pegou: " + livro.getTitulo() + " emprestado.");
        } else {
            System.out.println(
                    "[Empréstimo] Aviso: O cliente já possui o livro: "
                    + livro.getTitulo());
        }
    }

    public void listarEmprestados() {
        System.out.println("\n--- Livros de " + getNome() + " ---");

        // PERCORRENDO A LISTA
        // for(Tipo variavel : colecao) é o for-each do Java
        // equivale ao for x in lista do Python ou ao .each do Ruby
        for (Livros livro : livrosEmprestados) {
            System.out.println("- " + livro.getTitulo());
        }
    }
}
\end{lstlisting}

\begin{lstlisting}[style=java, caption={Exemplo de Aplicação de POO em Java: \texttt{Bibliotecario.java}}]
package org.example;

// HERANÇA + INTERFACE
// Bibliotecario herda de Pessoa e implementa GerenciaInventario
// por isso é obrigado a ter o método localizarLivro
public class Bibliotecario extends Pessoa implements GerenciaInventario {

    public Bibliotecario(String nome, int id) {
        super(nome, id);
    }

    // SOBRESCRITA (override)
    // @Override avisa o compilador que estamos implementando um método da
    // interface
    // equivale ao override do C# ou ao def com mesmo nome no Ruby
    @Override
    public Livros localizarLivro(String titulo, Estante estante) {
        System.out.println("\n[Busca] Bibliotecário " + getNome()
                + " localizando '" + titulo + "'...");

        // .values() retorna só os valores do HashMap, ignorando as chaves
        // aqui percorremos todos os Livros guardados na estante
        for (Livros livro : estante.getInventario().values()) {
            // equalsIgnoreCase compara strings ignorando maiúsculas e
            // minúsculas
            // equivale ao .downcase == do Ruby ou ao .lower() == do Python
            if (livro.getTitulo().equalsIgnoreCase(titulo)) {
                return livro;
            }
        }
        return null;
    }
}
\end{lstlisting}

\begin{lstlisting}[style=java, caption={Exemplo de Aplicação de POO em Java: \texttt{Main.java}}]
package org.example;

public class Main {
    public static void main(String[] args) {

        // INSTANCIAÇÃO
        // criando os objetos que vão ser usados ao longo da demonstração
        Autor billCipher    = new Autor("Bill Cipher");
        Autor kazuo         = new Autor("Kazuo Ishiguro");
        Autor williamGibson = new Autor("William Gibson");

        // HERANÇA em ação — Cliente e Bibliotecario herdam de Pessoa
        // (os dois SÃO Pessoas)
        Cliente gerard    = new Cliente("Gerard Way", 101);
        Bibliotecario oli = new Bibliotecario("Oliver Sykes", 201);

        // AGREGAÇÃO — a Estante é criada aqui e depois recebe Livros que
        // já existem
        Estante ficcao = new Estante("FIC-A1");


        // AGREGAÇÃO — Livro agrega um Autor que já existia antes
        Livros billsBook     = new Livros("978-0000000001",
                "Bill's Book", billCipher);
        Livros naoMeAbandone = new Livros("978-8535911732",
                "Não Me Abandone Jamais", kazuo);
        Livros neuromancer   = new Livros("978-8576572022",
                "Neuromancer", williamGibson);

        // ArrayList — palavras-chave mantém a ordem de inserção
        billsBook.adicionarPalavraChave("Misterio");
        billsBook.adicionarPalavraChave("Sobrenatural");
        naoMeAbandone.adicionarPalavraChave("Distopia");
        neuromancer.adicionarPalavraChave("Ficção Científica");
        neuromancer.adicionarPalavraChave("Cyberpunk");


        // HashMap — adicionando livros à estante mapeados pelo ISBN
        ficcao.adicionarLivro(billsBook);
        ficcao.adicionarLivro(naoMeAbandone);
        ficcao.adicionarLivro(neuromancer);

        // tentativa de título duplicado — o ArrayList com .contains() impede
        williamGibson.adicionarTitulo("Neuromancer");


        // AGREGAÇÃO — Cliente guarda referências pra Livros que já existiam
        gerard.emprestarLivro(billsBook);
        gerard.emprestarLivro(neuromancer);
        // tentativa de duplicata — o .contains() impede
        gerard.emprestarLivro(billsBook);
        gerard.listarEmprestados();


        // INTERFACE em ação
        // Bibliotecario implementa GerenciaInventario, então pode chamar
        // localizarLivro()
        String tituloBuscado = "Neuromancer";
        Livros encontrado = oli.localizarLivro(tituloBuscado, ficcao);

        if (encontrado != null) {
            System.out.println("-> Localizado com sucesso! Detalhes: "
                    + encontrado.getDetalhes());
        } else {
            System.out.println("-> Livro não encontrado.");
        }
    }
}
\end{lstlisting}

\section{Exercícios Práticos do Módulo}

Em cada exercício, todas as classes podem ficar em um único arquivo \texttt{.java}, como
explicado na seção ``Um Arquivo por Exercício'': a classe com o \texttt{main} vem primeiro e é a
única com \texttt{public}.

\begin{exercicio}{Exercício 1: Classe Abstrata e Herança (Os Moradores do Castelo)}
Em \textit{O Castelo Animado}, filme do Studio Ghibli, um castelo mágico anda sozinho pelos
campos, movido por um \textbf{demônio do fogo} que vive na lareira. Lá moram \textbf{magos} e o
próprio demônio, e todo morador sabe se apresentar, cada um do seu jeito. Não existe um
``morador genérico'': todo morador é de um tipo específico.

\textbf{Glossário:}
\begin{itemize}
    \item \textbf{Demônio do fogo}: criatura mágica feita de fogo, que vive na lareira e faz o
    castelo andar.
    \item \textbf{Feitiço}: magia que um mago sabe lançar. No código, é só um texto.
\end{itemize}

Implemente:

\begin{itemize}
    \item \texttt{Morador} (classe abstrata): atributo \texttt{private String nome}, recebido no
    construtor; getter \texttt{public String getNome()}; e método abstrato
    \texttt{public abstract void apresentar();}.

    \item \texttt{Mago}: herda de \texttt{Morador} e tem o atributo
    \texttt{private String feitico}. O construtor \texttt{Mago(String nome, String feitico)}
    repassa o nome com \texttt{super(nome)}. Sobrescreve \texttt{apresentar()} com
    \texttt{@Override}, imprimindo o nome (use \texttt{getNome()}) e o feitiço.

    \item \texttt{DemonioDoFogo}: herda de \texttt{Morador}. O construtor
    \texttt{DemonioDoFogo(String nome)} repassa o nome com \texttt{super(nome)}. Sobrescreve
    \texttt{apresentar()} com \texttt{@Override}, imprimindo o nome e que ele faz o castelo
    andar.
\end{itemize}

No programa principal, crie um \texttt{Mago} e um \texttt{DemonioDoFogo}, à sua escolha,
guardando os dois em variáveis do tipo \texttt{Morador}, e chame \texttt{apresentar()} nos dois.
Por fim, tente criar um objeto \texttt{Morador} diretamente e observe o erro: o programa nem
chega a compilar. Depois, transforme essa linha em comentário (\texttt{//}) para o programa
voltar a rodar.
\end{exercicio}

\begin{exercicio}{Exercício 2: Herança, Sobrescrita e ArrayList (O Cardápio do Leblanc)}
Em \textit{Persona 5}, o Leblanc é um pequeno café em Tóquio, famoso pelo café coado na hora e
pelo curry da casa. O cardápio tem itens comuns, vendidos pelo preço cadastrado, e
\textbf{combos}, que saem com desconto.

\textbf{Glossário:}
\begin{itemize}
    \item \textbf{Combo}: item que junta mais de um produto (por exemplo, curry e café) e custa
    10\% menos que o preço cadastrado.
\end{itemize}

Implemente:

\begin{itemize}
    \item \texttt{ItemDoCardapio}: atributos \texttt{private String nome} e
    \texttt{private double preco}, recebidos no construtor; getters
    \texttt{public String getNome()} e \texttt{public double getPreco()}; e o método
    \texttt{public double precoFinal()}, que retorna o preço.

    \item \texttt{ComboLeblanc}: herda de \texttt{ItemDoCardapio}. O construtor
    \texttt{ComboLeblanc(String nome, double preco)} repassa os dois valores com
    \texttt{super(nome, preco)}. Sobrescreve \texttt{precoFinal()} com \texttt{@Override},
    retornando o preço com 10\% de desconto (\texttt{getPreco() * 0.9}).

    \item \texttt{CafeLeblanc}: atributo \texttt{private ArrayList<ItemDoCardapio> cardapio},
    criado vazio no construtor, que não recebe parâmetros; método \texttt{public void adicionar(ItemDoCardapio item)}, que
    adiciona um item já criado ao cardápio; e método \texttt{public void mostrarCardapio()},
    que percorre o cardápio com um for-each e imprime o nome e o \texttt{precoFinal()} de cada
    item.
\end{itemize}

No programa principal, crie dois \texttt{ItemDoCardapio} e um \texttt{ComboLeblanc}, com nomes e
preços à sua escolha. Crie o \texttt{CafeLeblanc}, adicione os três itens e chame
\texttt{mostrarCardapio()}.
\end{exercicio}

\begin{exercicio}{Exercício 3: Interface e HashMap (Avistamentos na Cidade)}
Em \textit{Dandadan}, dois estudantes do ensino médio descobrem que fantasmas e alienígenas
existem de verdade e passam a anotar, em um mapa, os locais da cidade onde cada ser sobrenatural
foi visto. Alienígenas e espíritos não têm nada em comum, a não ser uma coisa: todos se
\textbf{manifestam} quando alguém investiga o local.

\textbf{Glossário:}
\begin{itemize}
    \item \textbf{Youkai}: espírito do folclore japonês, como um fantasma ou uma assombração.
    \item \textbf{Alienígena}: ser de outro planeta.
    \item \textbf{Avistamento}: registro de um ser sobrenatural visto em um local.
\end{itemize}

Implemente:

\begin{itemize}
    \item \texttt{Sobrenatural} (interface): método \texttt{void manifestar();}.

    \item \texttt{Alienigena}: implementa \texttt{Sobrenatural}. Tem o atributo
    \texttt{private String nome}, recebido no construtor, e implementa
    \texttt{public void manifestar()} com \texttt{@Override}, imprimindo o nome e um
    comportamento de alienígena à sua escolha.

    \item \texttt{EspiritoYoukai}: também implementa \texttt{Sobrenatural}, do mesmo jeito:
    atributo \texttt{private String nome}, recebido no construtor, e
    \texttt{public void manifestar()} com um comportamento de espírito à sua escolha.

    \item \texttt{MapaDeAvistamentos}: atributo
    \texttt{private HashMap<String, Sobrenatural> avistamentos}, criado vazio no construtor, que
    não recebe parâmetros; a chave do mapa é o nome do local. O método
    \texttt{public void registrar(String local, Sobrenatural ser)} guarda o ser no mapa, usando
    o local como chave. O método \texttt{public void investigar(String local)} busca o local no
    mapa: se não encontrar nada (\texttt{null}), imprime que o local está tranquilo; se
    encontrar, chama \texttt{manifestar()} do ser.
\end{itemize}

No programa principal, crie um \texttt{Alienigena} e um \texttt{EspiritoYoukai}, à sua escolha,
e registre cada um em um local à sua escolha. Depois, investigue os dois locais e um terceiro
local, que não foi registrado.
\end{exercicio}

\begin{exercicio}{Exercício 4: Revisão Geral (O Banquete da Masmorra)}
Em \textit{Dungeon Meshi} (\textit{Delicious in Dungeon}), um grupo de aventureiros explora uma
\textbf{masmorra} e, sem comida, aprende com um cozinheiro anão a preparar os monstros que
derrota. Todo monstro vira ingrediente, mas cada tipo tem a sua própria \textbf{receita}. O
cozinheiro carrega sempre a sua \textbf{panela}, que é dele e nasce com ele (composição), e
junta os monstros que o grupo derrota pelo caminho (agregação).

\textbf{Glossário:}
\begin{itemize}
    \item \textbf{Masmorra}: labirinto subterrâneo cheio de monstros.
    \item \textbf{Slime}: monstro gelatinoso, parecido com uma bolha.
    \item \textbf{Mandrágora}: planta-monstro com forma de gente, que grita quando é arrancada
    da terra.
    \item \textbf{Banquete}: refeição feita com todos os ingredientes reunidos.
\end{itemize}

Implemente:

\begin{itemize}
    \item \texttt{Monstro} (classe abstrata): atributo \texttt{private String nome}, recebido no
    construtor, e o getter \texttt{public String getNome()}; método abstrato
    \texttt{public abstract String receita();}, que nas filhas vai \textbf{retornar} o modo de
    preparo.

    \item \texttt{Slime} e \texttt{Mandragora}: herdam de \texttt{Monstro}. Cada uma tem um
    construtor (\texttt{Slime(String nome)} e \texttt{Mandragora(String nome)}) que repassa o
    nome com \texttt{super(nome)}, e sobrescreve
    \texttt{receita()} com \texttt{@Override}, retornando um modo de preparo à sua escolha.

    \item \texttt{Panela} (parte composta): atributo \texttt{private String material}, recebido
    no construtor, e o getter \texttt{public String getMaterial()}.

    \item \texttt{Cozinheiro}: atributos \texttt{private String nome},
    \texttt{private Panela panela} e \texttt{private ArrayList<Monstro> ingredientes}. O
    construtor \texttt{Cozinheiro(String nome, String materialDaPanela)} cria a \texttt{Panela}
    \textbf{dentro} dele, com o material recebido, e cria a lista de ingredientes vazia. O
    método \texttt{public void adicionarIngrediente(Monstro monstro)} adiciona um monstro já
    criado à lista. O método \texttt{public void servirBanquete()} imprime o nome do cozinheiro
    e o material da panela e, com um for-each, o nome e a \texttt{receita()} de cada
    ingrediente.
\end{itemize}

No programa principal, crie um \texttt{Slime} e uma \texttt{Mandragora}, à sua escolha,
\textbf{antes} do cozinheiro. Depois, crie o cozinheiro, com o nome e o material da panela à sua
escolha, adicione os dois monstros e chame \texttt{servirBanquete()}. Explique, em um comentário
no código, o que é composição e o que é agregação no seu programa.
\end{exercicio}
